%% toc: P15 161 161.2,161.4,161.5,161.9
% Durrett's remarks on normal approximation (Section 3.4) and the non-strict symmetric maximal inequality.
\chapter{Reading Normal Approximations: Notation, Corrections, Scales and Rates}
\label{sec:p15-w02c-p15-015-reading}

The central limit theorem of \cref{sec:p15-central-limit-theorems} is
a statement about limits, but it is used as a statement about finite
sums. Durrett surrounds his theorems in Section~3.4 with remarks that
say how the limit statements should be read and used: why the normal
variable carries the letter \(\chi\); what a roulette bet and a
histogram correction look like in a finite calculation; which
dependence hypothesis suffices for the strong law but not for the
normal limit; how a truncation level is chosen when the variance is
infinite and what happens when the tail is heavier still; at what size
of integer the prime-divisor theorem begins to say something; and how
sharp, and how necessary, the hypotheses of the Berry--Esseen bound are
\cite[Section~3.4, pp.~144--157]{DurrettPTEAuthor2019}. This chapter
states each of those remarks, proves every mathematical claim in them,
and corrects one stale cross-reference. It also proves in full the
non-strict case of the symmetric maximal inequality of
\cref{thm:p15-levy-symmetric-maximal}.

\emph{Inputs.} Finite real random variables on a probability space; the
notation of \cref{def:p15-clt-notation}; for the maximal inequality,
finitely many independent summands with laws symmetric about zero.
\emph{Conclusions.} The distribution of \(\chi^2\)
(\cref{thm:p15-w02c-p15-015-chi-square}); the exact consistency of the
histogram correction (\cref{thm:p15-w02c-p15-015-histogram}); the
separation of the pairwise strong law from the normal limit
(\cref{thm:p15-w02c-p15-015-pairwise-contrast}); the admissible
truncation window of the infinite-variance example
(\cref{thm:p15-w02c-p15-015-truncation-window}); a nonnormal limit for
tails \(x^{-\alpha}\), \(0<\alpha<2\)
(\cref{thm:p15-w02c-p15-015-power-tail-limit}); the summability that
an \(n^{-1/2}\) rate implies (\cref{thm:p15-w02c-p15-015-rate-series});
and the non-strict maximal inequality
(\cref{thm:p15-w02c-p15-015-levy-nonstrict}).
\emph{Standing assumptions.} All random variables are finite, real and
measurable; \(\Phi\), \(\chi\), \(S_n\) and \(\Rightarrow\) are as in
\cref{def:p15-clt-notation}. Numerical values written to four decimal
places are rounded computations of the stated exact expressions; they
are estimates, not identities, unless the text says ``exact''.

\section{The letter \texorpdfstring{\(\chi\)}{chi} for a standard normal variable}
\label{sec:p15-w02c-p15-015-chi}

Durrett writes \(\chi\) for a standard normal variable in his
Theorem~3.4.1 and comments at once that the choice departs from common
usage but is convenient: the reason is that the square of a normal
variable has a chi-squared distribution
\cite[p.~144]{DurrettPTEAuthor2019}. The letter therefore records the
link between the normal variable and the chi-squared law of its
square. The following lemma states that link exactly.

\begin{hblemma}{thm:p15-w02c-p15-015-chi-square}{proved=annex:p15-w02c-p15-015-chi-square}{The square of a standard normal variable}
Let \(\chi\) have the standard normal law. Then \(\chi^2\) has the
density
\[
 f(y)=(2\pi y)^{-1/2}e^{-y/2}\quad(y>0),\qquad f(y)=0\quad(y\leq0),
\]
which is the chi-squared density with one degree of freedom. In
particular \(\mathbb E\chi^2=1\). For \(\sigma\geq0\), the square of
\(\sigma\chi\) has the law of \(\sigma^2\chi^2\).
\end{hblemma}

\section{Roulette, the histogram correction, and a small sample}
\label{sec:p15-w02c-p15-015-histogram}

\begin{hbexample}{def:p15-w02c-p15-015-roulette}{The roulette model behind the gain law}
Durrett's Example~3.4.5 starts from a wheel with slots numbered
\(1\) to \(36\), eighteen of them red and eighteen black, and two
further green slots numbered \(0\) and \(00\)
\cite[p.~145]{DurrettPTEAuthor2019}. A player stakes one dollar on red
(or on black); if the ball stops in a slot of that colour the player
wins one dollar, and otherwise loses the stake. With the ball equally
likely to stop in each of the \(38\) slots and the spins independent,
the gain \(X_i\) on spin \(i\) takes the value \(1\) with probability
\(18/38\) (the eighteen slots of the chosen colour) and \(-1\) with
probability \(20/38\) (the other colour and the two green slots). This
is the gain law of \cref{thm:p15-clt-numerical-examples}(1): mean
\(-1/19\), variance \(1-(1/19)^2=360/361\). The two green slots are the
whole of the house advantage: without them the bet would be fair.
\end{hbexample}

Durrett's Example~3.4.7 approximates \(\mathbb P(S_{16}=8)\) for the
number of heads in sixteen fair tosses by treating the integer \(8\) as
the interval \([7.5,8.5]\), and obtains \(0.1974\); he observes that
this is close to the exact value \(\binom{16}8 2^{-16}=0.1964\), even
though \(n=16\) is small \cite[p.~146]{DurrettPTEAuthor2019}. He then
names the device the histogram correction, says it matters when the
normal approximation is used with small \(n\), and gives a rule and a
reason: \(\mathbb P(S_{16}\leq11)\) is to be read as
\(\mathbb P(S_{16}\leq11.5)\), and one motive is that this gives the
same answer as reading the same probability as
\(1-\mathbb P(S_{16}\geq12)\). The next result shows that this
consistency holds exactly, for every integer threshold and every
continuous approximating distribution, and fails without the
correction.

\begin{hbmethod}{def:p15-w02c-p15-015-histogram-rule}{The histogram correction}
Let \(S\) take only integer values and let \(G\) be a continuous CDF
used to approximate the CDF of \(S\); for a sum with mean \(a\) and
standard deviation \(b>0\) the normal choice is
\(G(x)=\Phi((x-a)/b)\). For integers \(j\leq k\), the corrected
approximations are
\[
 \mathbb P(S\leq k)\approx G(k+\tfrac12),\qquad
 \mathbb P(S\geq k)\approx1-G(k-\tfrac12),\qquad
 \mathbb P(j\leq S\leq k)\approx G(k+\tfrac12)-G(j-\tfrac12).
\]
Each integer \(m\) is represented by the interval
\([m-\tfrac12,m+\tfrac12]\), and an event about \(S\) by the union of
the intervals of its integers. Side conditions: \(S\) must be
integer-valued (for a lattice of span \(h\) and offset \(c\), apply the
rule to \((S-c)/h\)); \(G\) must be continuous so that the endpoints of
the intervals carry no mass; and the result is an approximation whose
error the rule itself does not bound (\cref{thm:p15-clt-berry-esseen}
gives a bound for normalized sums with a finite third moment).
\end{hbmethod}

\begin{hbtheorem}{thm:p15-w02c-p15-015-histogram}{proved=annex:p15-w02c-p15-015-histogram}{Exact consistency of the histogram correction}
Let \(S\) be integer-valued and \(G\) a continuous CDF. For every
integer \(k\), the corrected approximations of
\(\mathbb P(S\leq k)\) and \(\mathbb P(S\geq k+1)\) sum to exactly
\(1\), as the probabilities themselves do; and the corrected
approximations of \(\mathbb P(S=m)\), \(m\in\mathbb Z\), sum to
exactly \(1\). Without the correction, the approximations
\(G(k)\) of \(\mathbb P(S\leq k)\) and \(1-G(k+1)\) of
\(\mathbb P(S\geq k+1)\) sum to \(1-(G(k+1)-G(k))\). For a normal
approximation \(G(x)=\Phi((x-a)/b)\) with \(b>0\), this sum is strictly
less than \(1\) for every \(k\), so the two uncorrected readings of
the same probability always disagree.
\end{hbtheorem}

\begin{hbexample}{def:p15-w02c-p15-015-sixteen-tosses}{Sixteen tosses, with and without the correction}
Let \(S_{16}\) be the number of heads in sixteen fair tosses, so
\(\mathbb ES_{16}=8\), \(\operatorname{Var}S_{16}=4\), and
\(G(x)=\Phi((x-8)/2)\).
\begin{enumerate}
\item \emph{A point probability.} The corrected value is
 \(G(8.5)-G(7.5)=2\Phi(1/4)-1\approx0.1974\). The exact value is
 \(\binom{16}8 2^{-16}=12870/65536\approx0.1964\). Durrett's comparison
 of these two numbers is the evidence he gives that the correction works
 even for \(n=16\) \cite[p.~146]{DurrettPTEAuthor2019}.
\item \emph{A cumulative probability.} Exactly,
 \(\mathbb P(S_{16}\geq12)=\bigl(\binom{16}{12}+\binom{16}{13}+\binom{16}{14}
 +\binom{16}{15}+\binom{16}{16}\bigr)2^{-16}=2517/65536\), so
 \(\mathbb P(S_{16}\leq11)=63019/65536\approx0.9616\). The corrected
 reading \(\mathbb P(S_{16}\leq11.5)\) gives \(\Phi(1.75)\approx0.9599\),
 and the corrected reading of \(1-\mathbb P(S_{16}\geq12)\) gives
 \(1-(1-\Phi(1.75))\), the same number. Uncorrected, the first reading
 gives \(\Phi(1.5)\approx0.9332\) and the second gives
 \(\Phi(2)\approx0.9772\): two different answers to one question, both
 further from the exact value than the corrected one.
\end{enumerate}
\end{hbexample}

\section{Pairwise independence: enough for the strong law, not for the normal limit}
\label{sec:p15-w02c-p15-015-pairwise}

Durrett opens his Example~3.4.9 with a contrast: pairwise independence
is a sufficient hypothesis for his strong law of large numbers,
Theorem~2.4.1 \cite[p.~76]{DurrettPTEAuthor2019}, but it is not
sufficient for the central limit theorem
\cite[p.~147]{DurrettPTEAuthor2019}. The two halves are
\cref{thm:p15-etemadi-strong-law} and
\cref{thm:p15-clt-pairwise-boundary}. The next corollary puts them
side by side for one sequence.

\begin{hbcorollary}{thm:p15-w02c-p15-015-pairwise-contrast}{proved=annex:p15-w02c-p15-015-pairwise-contrast}{One pairwise independent sequence, two verdicts}
Let \((X_j)\) be the pairwise independent, identically distributed
sequence of fair signs of \cref{thm:p15-clt-pairwise-boundary}, and
\(S_n=\sum_{j=1}^nX_j\). Then \(S_n/n\to0=\mathbb EX_1\) almost surely,
and \(\operatorname{Var}(S_n)=n\); but \(S_n/\sqrt n\) does not
converge weakly to \(\chi\), nor to any normal law with positive
variance. The failure is visible along \(n=2^r\), where
\(\mathbb P(|S_{2^r}|/2^{r/2}\leq\frac12)\geq1-2^{-r}\to1\), while
\(\mathbb P(|\chi|\leq\frac12)=2\Phi(\frac12)-1<0.4\).
\end{hbcorollary}

The distinction is one of order. The strong law concerns averages
\(S_n/n\), and Etemadi's argument needs only the variances of truncated
sums, which pairwise independence makes additive. The normal limit
concerns the whole law of \(S_n/\sqrt n\), which depends on the joint
behaviour of more than two summands at a time; pairwise independence
says nothing about that, and the construction exploits exactly this.

\section{Choosing a truncation level, and tails heavier than the inverse square}
\label{sec:p15-w02c-p15-015-truncation}

In Durrett's Example~3.4.13 the summands are symmetric with
\(\mathbb P(|X_1|>x)=x^{-2}\) for \(x\geq1\), and the truncation level is
\(c_n=\sqrt n\log\log n\) (\cref{thm:p15-clt-infinite-variance-example}).
His stated reason for this choice has two sides
\cite[p.~151]{DurrettPTEAuthor2019}. The level must be high enough that,
with probability tending to one, no summand is changed by the truncation.
Since the variance of the truncated summands should be as small as
possible, it should also stay near the lowest level that achieves this.
The next lemma makes both sides precise and shows what goes wrong at the
two extremes.

\begin{hblemma}{thm:p15-w02c-p15-015-truncation-window}{proved=annex:p15-w02c-p15-015-truncation-window}{The admissible truncation window for the inverse-square tail}
Let \(X_1,X_2,\ldots\) be i.i.d., symmetric, with
\(\mathbb P(|X_1|>x)=x^{-2}\) for \(x\geq1\), and let \(c_n\geq1\) be
deterministic. Put \(Y_{n,m}=X_m\mathbf1_{\{|X_m|\leq c_n\}}\),
\(1\leq m\leq n\).
\begin{enumerate}
\item \(\mathbb P(Y_{n,m}\ne X_m\text{ for some }m\leq n)\leq n/c_n^2\),
 and \(n\mathbb P(|X_1|>c_n)=n/c_n^2\). This tends to zero if and only
 if \(c_n/\sqrt n\to\infty\).
\item \(\mathbb EY_{n,1}=0\) and \(\mathbb EY_{n,1}^2=2\log c_n\),
 which increases with \(c_n\). Moreover
 \(\mathbb EY_{n,1}^2\sim\log n\) if and only if
 \(\log(c_n/\sqrt n)=o(\log n)\); a level \(c_n=n^{1/2+\varepsilon}\)
 with \(\varepsilon>0\) gives \((1+2\varepsilon)\log n\).
\item If \(c_n/\sqrt{n\log n}\to0\), then for every \(\varepsilon>0\)
 the Lindeberg sum of the row \(Y_{n,m}/\sqrt{n\log n}\) is zero for all
 large \(n\).
\item The source level \(c_n=\sqrt n\log\log n\) (for \(n\geq16\), where
 \(\log\log n>1\)) satisfies the conditions of items 1--3:
 \(n/c_n^2=(\log\log n)^{-2}\to0\),
 \(\log(c_n/\sqrt n)=\log\log\log n=o(\log n)\), and
 \(c_n/\sqrt{n\log n}=\log\log n/\sqrt{\log n}\to0\).
\item At the other extreme, \(c_n=n\) still satisfies item~1, but
 \(\mathbb EY_{n,1}^2=2\log n\), twice the source value; and the
 condition of item~3 fails, since \(c_n/\sqrt{n\log n}\to\infty\).
\end{enumerate}
\end{hblemma}

Item~1 forces \(c_n\) above \(\sqrt n\); item~2 shows that every
level allowed by item~1 keeps \(\mathbb EY_{n,1}^2\) above \(\log n\)
for large \(n\), and that a level of order \(n^{1/2+\varepsilon}\)
inflates it by the factor \(1+2\varepsilon\); item~3 is what makes the Lindeberg condition
of the source's proof hold trivially. The level \(\sqrt n\log\log n\)
sits just above the lower edge of the window, which is Durrett's
``close to the lowest possible level''.

Durrett follows the example with a remark: if the inverse-square tail is
replaced by \(\mathbb P(|X_1|>x)=x^{-\alpha}\) with \(0<\alpha<2\), the
sum normalized by \(n^{1/\alpha}\) converges to a limit that is not
\(\chi\) \cite[p.~152]{DurrettPTEAuthor2019}. The source sends the
reader to ``Section~3.6'' for this. In this version of the text
Section~3.6 is Poisson convergence; stable laws are Section~3.8, which
the contents place at p.~182 \cite[p.~182]{DurrettPTEAuthor2019}. The
remark is proved directly here.

\begin{hbtheorem}{thm:p15-w02c-p15-015-power-tail-limit}{proved=annex:p15-w02c-p15-015-power-tail-limit}{Power tails below the inverse square give a nonnormal limit}
Let \(0<\alpha<2\), and let \(X_1,X_2,\ldots\) be i.i.d. with laws
symmetric about zero and \(\mathbb P(|X_1|>x)=x^{-\alpha}\) for
\(x\geq1\). Put
\[
 I_\alpha=\int_0^\infty(1-\cos y)\,\alpha y^{-\alpha-1}\,dy .
\]
Then \(0<I_\alpha<\infty\), and
\(S_n/n^{1/\alpha}\Rightarrow Q_\alpha\), where \(Q_\alpha\) is the
probability law with characteristic function
\(\exp(-I_\alpha|t|^\alpha)\). The law \(Q_\alpha\) is not normal: it is
neither the law of \(\chi\), nor \(N(m,s^2)\) for any \(m\in\mathbb R\)
and \(s\geq0\), the degenerate case included. No centering is needed,
for every \(\alpha\in(0,2)\), including \(0<\alpha\leq1\) where
\(\mathbb E|X_1|=\infty\).
\end{hbtheorem}

For \(\alpha=2\) the integral defining \(I_2\) diverges near zero, so
the computation gives no finite limit at the scale \(\sqrt n\): the
inverse-square tail needs the extra factor \(\sqrt{\log n}\) of
\cref{thm:p15-clt-infinite-variance-example}, and then the limit is
normal. For \(\alpha<2\) the scale \(n^{1/\alpha}\) grows faster than
\(\sqrt n\).

\section{L\'evy's normal-convergence criterion and its source}
\label{sec:p15-w02c-p15-015-levy-criterion}

Durrett introduces his Theorem~3.4.14, the necessary and sufficient
condition for \((S_n-a_n)/b_n\Rightarrow\chi\) with suitable constants,
as the final answer to the question of convergence to the normal law,
and attributes it to L\'evy \cite[p.~152]{DurrettPTEAuthor2019}. He
gives no proof. He refers instead to Gnedenko and Kolmogorov (1954), and
describes that monograph as the definitive account of many questions
about sums of independent variables. In this handbook the criterion, both its
directions, and the construction of the normalizers are proved in
\cref{thm:p15-nda-characterization,thm:p15-nda-normalizer-construction},
and the cutoff route of the 1954 book is
\cref{thm:p15-nda-gk-cutoff-route}. The infinite-variance example of
\cref{thm:p15-clt-infinite-variance-example} satisfies the criterion,
and the power tails of \cref{thm:p15-w02c-p15-015-power-tail-limit} do
not, since for them
\(y^2\mathbb P(|X_1|>y)/\mathbb E(X_1^2;|X_1|\leq y)\to(2-\alpha)/\alpha\ne0\)
(\cref{thm:p15-nda-examples}, item~3).

\section{How large an integer the prime-divisor theorem needs}
\label{sec:p15-w02c-p15-015-erdos-kac-scale}

The Erd\H{o}s--Kac theorem (\cref{thm:p15-clt-erdos-kac}) says that a
uniformly chosen integer up to \(n\) has about
\(\log\log n+\chi(\log\log n)^{1/2}\) distinct prime divisors. Durrett
opens his Section~3.4.3 by noting that \(\exp(e^4)\) is about
\(5.15\times10^{23}\), and concludes that the result does not describe
most of the numbers met in everyday life
\cite[p.~153]{DurrettPTEAuthor2019}. The arithmetic behind the remark is
this. Since \(e^4\approx54.598\), the centering \(\log\log n\) reaches
\(4\), and the spread \((\log\log n)^{1/2}\) reaches \(2\), only when
\(n\geq\exp(e^4)\approx5.15\times10^{23}\). For \(n=10^9\),
\(\log\log n\approx3.03\) and the spread is about \(1.74\); for
\(n=10^{12}\), \(\log\log n\approx3.32\). On such ranges the centering
and the spread are only a few units, so the normal law is being compared
with a count that takes only a handful of likely values. The theorem
describes the regime in which \(\log\log n\) is large, and that regime
begins far beyond ordinary magnitudes.

\section{The Berry--Esseen constant, its source, and the third moment}
\label{sec:p15-w02c-p15-015-berry-esseen-remarks}

After his Theorem~3.4.17 (\cref{thm:p15-clt-berry-esseen}), Durrett adds
several remarks \cite[p.~157]{DurrettPTEAuthor2019}.
\begin{enumerate}
\item The constant \(3\) is not the best known; he reports that van Beek
 (1972) obtains \(0.8\).
\item He endorses the method by quoting Feller, whose own assessment of
 the smoothing proof reads: ``Our streamlined method yields a remarkably
 good bound even though it avoids the usual messy numerical
 calculations'' \cite[ch.~XVI, sec.~5, p.~542, n.~7]{Feller1971ProbabilityII}.
 In the same note Feller lists earlier constants and states that no
 substantial improvement can be expected without improving the error
 term of his smoothing inequality, which is the term
 \(24\lambda/(\pi L)\) of \cref{thm:p15-clt-cdf-inversion}.
\item He states that the third-moment hypothesis is what is needed for
 the rate \(n^{-1/2}\).
\item He reports Heyde's (1967) theorem: for \(0<\delta<1\), with
 \(F_n\) and \(\Phi\) as in \cref{thm:p15-clt-berry-esseen},
 \[
  \sum_{n=1}^\infty n^{-1+\delta/2}\sup_x|F_n(x)-\Phi(x)|<\infty
  \quad\text{if and only if}\quad\mathbb E|X_1|^{2+\delta}<\infty .
 \]
\item For this and further results on rates he refers to Hall (1982).
\end{enumerate}
The third remark should be read together with the fourth. The following
lemma gives the one step that connects them and that can be proved
from the definitions alone.

\begin{hblemma}{thm:p15-w02c-p15-015-rate-series}{proved=annex:p15-w02c-p15-015-rate-series}{An inverse-square-root rate makes Heyde's series converge}
Let \(\Delta_n\geq0\) satisfy \(\Delta_n\leq Cn^{-1/2}\) for all
\(n\geq1\) and some finite \(C\). Then
\(\sum_{n\geq1}n^{-1+\delta/2}\Delta_n<\infty\) for every
\(0<\delta<1\). If instead \(\sum_nn^{-1+\delta/2}\Delta_n=\infty\) for
some \(0<\delta<1\), then \(\sqrt n\,\Delta_n\) is unbounded.
\end{hblemma}

Applied to \(\Delta_n=\sup_x|F_n(x)-\Phi(x)|\), the lemma and Heyde's
theorem as Durrett reports it together say: an \(O(n^{-1/2})\) rate
forces \(\mathbb E|X_1|^{2+\delta}<\infty\) for every \(\delta<1\), and
a law with \(\mathbb E|X_1|^{2+\delta}=\infty\) for some \(\delta<1\)
has no \(O(n^{-1/2})\) rate. This is the precise sense in which the
moment hypothesis of order three cannot be lowered to any order
\(2+\delta<3\) in \cref{thm:p15-clt-berry-esseen}. Durrett does not
show, and nothing here shows, that finiteness of the third moment itself
is necessary for an \(O(n^{-1/2})\) rate.

\begin{proofboundary}{def:p15-w02c-p15-015-reported-results}{Results reported by the source without proof}
Three results mentioned in this chapter are attributions that Durrett
reports and does not prove: van Beek's constant \(0.8\), Heyde's
equivalence, and, for the extreme value limits, Gnedenko's (1943)
conditions recorded after \cref{thm:p15-wce-extreme-maxima}; Hall (1982)
is a reading pointer. They are stated here as the source states them, with the source's
locator, and no conclusion of this chapter's numbered results depends on
them. L\'evy's criterion, which Durrett also attributes without proof,
is proved in \cref{thm:p15-nda-characterization}. The general theory
of stable laws, to which Durrett defers the power-tail remark, is a
separate presentation; the remark itself is proved above.
\end{proofboundary}

\section{The non-strict symmetric maximal inequality}
\label{sec:p15-w02c-p15-015-levy-nonstrict}

\Cref{thm:p15-levy-symmetric-maximal} states the maximal inequality for
partial sums of independent symmetric summands with strict thresholds,
and adds that it also holds with non-strict thresholds. The non-strict
form is the one needed when the event \(\{|T_n|\geq t\}\) rather than
\(\{|T_n|>t\}\) is controlled, as in the non-strict comparison of
\cref{thm:p15-feller-largest-term-inequality}.

\begin{hblemma}{thm:p15-w02c-p15-015-levy-nonstrict}{proved=annex:p15-w02c-p15-015-levy-nonstrict}{The non-strict case of the symmetric maximal inequality}
Let \(Y_1,\ldots,Y_n\) be independent real random variables whose
individual laws are symmetric about zero, and put \(T_0=0\),
\(T_k=\sum_{i=1}^kY_i\). For every \(t>0\),
\[
 \mathbb P\!\left(\max_{1\leq k\leq n}|T_k|\geq t\right)
          \leq2\,\mathbb P(|T_n|\geq t).
\]
\end{hblemma}

The constant \(2\) cannot be lowered in general: for \(n=2\) and
independent fair signs \(Y_1,Y_2\), with \(t=1\), the left side is
\(1\) and \(\mathbb P(|T_2|\geq1)=1/2\). Symmetry cannot be dropped: for
\(n=2\), \(Y_1=1\) and \(Y_2=-1\) (constants, not symmetric), and
\(t=1\), the left side is \(1\) while \(\mathbb P(|T_2|\geq1)=0\).

\section{Relationships and sources}
\label{sec:p15-w02c-p15-015-relationships}

The results above refine the main chapter on central limit theorems,
\cref{sec:p15-central-limit-theorems}: the notation of
\cref{def:p15-clt-notation}, the numerical examples of
\cref{thm:p15-clt-numerical-examples}, the pairwise construction of
\cref{thm:p15-clt-pairwise-boundary}, the infinite-variance example of
\cref{thm:p15-clt-infinite-variance-example}, the Erd\H{o}s--Kac theorem
\cref{thm:p15-clt-erdos-kac}, and the Berry--Esseen bound
\cref{thm:p15-clt-berry-esseen}. The power-tail theorem uses the product
lemma \cref{thm:p15-clt-complex-products} and the continuity theorem
\cref{thm:p15-characteristic-continuity-theorem}. The non-strict maximal
inequality completes \cref{thm:p15-levy-symmetric-maximal} in
\cref{sec:p15-weak-laws-and-triangular-arrays}. The concluding clause of
Durrett's Exercise~3.3.24, that a finite essential bound gives the
moment-method hypothesis, is proved in
\cref{thm:p15-w02c-p15-012-bound-gives-growth}; the weak-law chapter's
pointer to the Section~3.3 exercises relies on it.

Sources: Durrett, author Version~5 dated 11 January 2019,
Section~3.4 \cite[pp.~144--157]{DurrettPTEAuthor2019}, Theorem~2.4.1
\cite[p.~76]{DurrettPTEAuthor2019}, Exercise~3.2.2
\cite[p.~123]{DurrettPTEAuthor2019} and the contents
\cite[p.~182]{DurrettPTEAuthor2019}; Feller's second volume, second
edition, for the quoted assessment of the smoothing method
\cite[ch.~XVI, sec.~5, p.~542, n.~7]{Feller1971ProbabilityII}.
